\paragraph{Outward verification of every CSR grid menu.}
A further verification study certifies all 12 CSR confirmation menus and
both CSR grid scaling menus, through $n=4096$, using interval $LDL^T$
elimination in band storage. All 14 certificates pass standalone replay,
with global gaps below $1.01\times10^{-5}$ (Table~\ref{tab:banded}).
This strengthens the preceding numerical checks without changing any
optimizer, saved menu, or isolated runtime measurement.

An initial binary64 interval implementation rejected a rank-eight witness
because its eigenvalue brackets became too wide to prove exclusion. That
attempt is retained. The revised verifier retries unresolved counts at
$n\le512$ with directed 50-digit decimal intervals; 57 counts use this
path. Both larger grid cases need only outward binary64 arithmetic.
The final checker independently encloses ordered eigenvalue sums and actual
polar-orthonormalized projectors; it imports no menu solver. Six CPU
processes check independent price/site pairs. These parallel verification
times are separate from the isolated optimization study.

For half-bandwidth $b$, each unpivoted inertia count uses $O(nb^2)$ work
and $O(nb+b^2)$ storage. The grid has $b=\lfloor\sqrt n\rfloor$;
at $n=4096$, packed input bands occupy about 6.1 MiB and an interval factor
about 4.1 MiB per worker, excluding temporary arrays and process overhead.
This is an advantage of band structure, not a guarantee for arbitrary
sparse graphs. Ambiguous pivots cause rejection; high bandwidth can erase
the storage benefit. Appendix~\ref{app:sturm} gives the factorization
argument and arithmetic assumptions.

\begin{table}[t]\centering\small
\caption{Independently replayed outward certificates for CSR grid menus.
Decimal counts record successful higher-precision inertia retries;
the displayed gap column is rounded. All unrounded gaps are below
$1.01\times10^{-5}$.}
\label{tab:banded}\begin{tabular}{rrrrr}\toprule
$n$ & Menus & Half-bandwidth & Decimal counts & Max. gap / $10^{-5}$\\\midrule
128 & 4 & 11 & 18 & 1.000029 \\
256 & 4 & 16 & 26 & 1.000031 \\
512 & 4 & 22 & 13 & 1.000031 \\
1024 & 1 & 32 & 0 & 1.000008 \\
4096 & 1 & 64 & 0 & 1.000008 \\
\bottomrule\end{tabular}

\end{table}
